% 来源：OEBPS/Text/Section0001_split_004.xhtml#sigil_toc_id_35

\BookChapter{24}
\section[第24章\q 梦想水桶]{第24章\\ 梦想水桶}

\begin{BookQuote}
停止了梦想，就停止了生命。

\hfill ——马尔科姆·福布斯
\end{BookQuote}

什么是梦想水桶？这是一个特别的地方，你特意把有些东西放在这里，只是为了你自己还有那些你非常热爱的人可以享受人生，让你一边创造财富一边享受财富创造的梦想生活。这些东西是为了现在，而不是为了明天！你的梦想水桶意味着要让你兴奋，给你的生活注入甜蜜，让你愿意赚得更多，奉献得更多。想象一下，哪些东西你愿意攒钱去买，哪些东西你愿意攒钱去做，这些东西就可以放进你的梦想水桶作为一种战略性挥霍物资，其实就是有计划、有目的的败家。

现在让你的船开起来吧。也许，你会给自己买双马诺洛·布拉尼克高跟鞋——你想了好久了，或者买一个迈阿密热火队的场边座位，或者带着孩子到迪士尼享受一次贵宾之旅，或者你开始往梦想水桶装更多的钱，追求更大的梦想：赛季套票。你也许夏天会进行一次登山之旅，冬天去滑冰或者单板滑雪，好好度个假。你也许会买辆新车，不是实用型，而是像迷你库柏那种时髦型的或者野马汽车那种彪悍型的。你也许会买座度假公寓或者度假别墅。

我认识一个百万富翁，坐飞机总是坐经济舱，因为他总想多省些钱，但是他的妻子经常抱怨：“我们有足够的钱，为什么不能享受一下？”这经常会成为他们夫妻之间冲突的导火索，因为他们为了生意上的事经常要坐飞机出差。参加了我的投资理财大师训练营之后，他决定用他梦想水桶里的钱来享受一下生活，他和他的家人外出旅行的时候会升级到公务舱。结果他发现，只是多花一些钱，升级一下舱位，不但让他自己的旅行变得更加舒服了，而且让他的家庭生活更加和谐了（这事其实更重要）。也许有一天，他会考虑包架私人飞机，民航飞机也不坐了——坐私人飞机也许并没有他想象的那么贵。

很多人有了很多钱后，他们的生活方式还跟以前一样，没有品质可言。他们一天天看着银行账户上的财富数字积累得越来越大，却错过了生活的快乐和幸福。本来在创造财富的同时，他们也可以自己创造快乐和幸福，并分享快乐和幸福。

我第一次想到梦想水桶这个主意，是在我的事业发展初期。那个时候我能想到的头奖是能买上两套新的西装，因为品牌男装连锁商店正在打折，或者到夏威夷度个假。当时对于我来说，这可是要花一大笔钱的。

我在斐济的那个度假村，是一个更大的梦想，后来梦想成真了。正如我跟你分享过的那样，那是在我24岁的时候，我爱上了南太平洋岛上那绿松石般的海水。我的心就像一下子找到了家。我一直想给自己、我的朋友、我的家人找到一个世外桃源。现在，经过了这么多年的努力之后，马雷度假村和水疗中心变成了规模相当大的资产，因为我不断地打造它，把它变成了南太平洋上的顶级旅游度假胜地。但是这只是一个“红包”，事实上，马雷度假村十多年来一直是斐济排名第一的度假胜地。2013年，奥普拉选择这里作为她最喜欢度假的地方。这本来只是我梦想做的事情，结果梦想成真，我创造出来这样一个美好的度假村。

你的梦想不是用来给你赚钱发财的，而是用来给你带来更好生活品质的。这不正是你首先要填满前面两个水桶的原因吗？你必须先练习自我约束，你在向梦想水桶里放钱的时候，也必须约束自我，和前面两个水桶一样。如果你把所有的钱都放到梦想水桶里，你可能最终会破产，就像摇滚巨星威利·尼尔森一样，所以关键在于平衡。你梦想水桶中的奖金，不一定只是为了你自己。最好的是那些你给予他人的东西。

\begin{BookQuote}
梦想是你性格的试金石。

\hfill ——亨利·戴维·梭罗
\end{BookQuote}

也许你跟我一样，就是喜欢给别人送礼物。最好的礼物就是那些对方意想不到的礼物。

我妈妈在年轻的时候没有钱，我们一家人总是挣扎着过日子，住在洛杉矶东部一个很便宜的破房子里，那里几乎天天都会播放雾霾警报，告诉我们在户外行走，呼吸空气不安全。

所以，后来在我生意开始好起来之后不久，一天，我让我妈妈陪我去看一套公寓，我想把它买下来，就在亨廷顿海滩的附近。我带她里里外外地看了一遍那套公寓，让她看看美妙的海景。然后我们从房子里走出来，来到了海滩上，呼吸着略带咸味的空气。

我说：“妈妈，我真的很喜欢这个地方。但是我希望你帮我做决定。一句话，你觉得这套房子怎么样？”

“你是开玩笑吧？”我妈妈说，“这根本不可能！你想一想我们是从什么地方出来的，现在你竟然想住在这里？”

“那么你觉得这个地方好吗，妈妈？”我问道。

“嗯，这个地方太好了，好得简直让人难以相信！”妈妈回答。

于是，我递给妈妈一串钥匙。

“这是什么？”妈妈问我。

“这套房子就是你的了。”我说。

我永远不会忘记妈妈脸上那震惊的表情以及快乐的泪水。我妈妈现在已经过世了。但是我还记得那个时刻，每个细节我都记得清清楚楚，那是我一生最喜欢回味的时刻。

你并不需要等待，你可以这样做。你可以实现你的梦想。如果你非常想要实现梦想，你肯定能找到办法。

我给我母亲买了那套公寓后不久，我接待了一个小学生团队，有100多个人，他们都是五年级的学生，来自邻近的贫困地区，得克萨斯州休斯敦市下面的一个小学。他们大多数人从小就形成了思维定式，认为自己绝对上不了大学，他们的生活习惯和学习习惯也不好，当然这样下去他们也绝对上不了大学。于是我当场决定，跟他们签订一份合同。我会给他们支付4年的大学学费，只要他们能够保持平均成绩至少达到B而且不惹麻烦就行。我说得很清楚，只要能够专注，任何一个人都能超过平均水平，我会提供指导来帮助他们。我和他们约法三章：第一，他们必须约束自己，别惹麻烦，不能进监狱。第二，在高中毕业之前，女生不能怀孕。第三，最重要的是，他们每年需要奉献出20个小时服务自己所在的社区，到一些慈善机构做义工。为什么我要加上第三条？上大学当然非常好，但是对于我来说，更加重要的是要教育他们认识到，自己可以给予别人一些东西，他们的人生并不是要从别人那里得到某些东西。我当时签订这个助学合同的时候，根本不知道长期而言我要帮多少人支付学费，但是我会百分之百地遵守承诺，而且我签订的是一个有法律约束力的合同，它要求我必须遵守合同提供这些助学金。这很有意思，你看，我别无选择，只能前进。我总是说，如果你想占领这个岛屿的话，你必须先烧掉你的船，像中国人说的那样，破釜沉舟！于是我和这100多个五年级的小学生签订了这样的合同。后来这些孩子中有23个和我一起合作，一起努力，从五年级开始一直到他们上了大学。后来还有几个孩子读了研究生，其中还有法学院的研究生。我称他们为“我的冠军”。现在这23个上了大学的孩子，有的成为社会工作者，有的创办企业当了老板，有的为人父母。就在几年前，我们重新聚到了一起。我听他们讲述自己那些神奇的故事，小时候，他们每年都做过20个小时的义工，这种能够给予别人的事情，给他们幼小的心灵留下了深刻的印象，后来这成了他们持续一生的固定模式。做义工义务奉献，让他们相信自己的人生是有价值、有意义的。他们现在把这一点教给他们自己的孩子。

我现在告诉你这个故事，是因为你不需要等到你完全准备好了才去实现你的梦想。你只要马上去做，就能找到路径，而且恩泽会来找你。你承诺去做一些事，更多的是去服务别人，而不仅是服务你自己——有些人称之为运气，有些人称之为巧合。至于你相信什么，还是由你自己来决定。你只要知道，当你给出了你的所有，奉献出了你的所有，福报会是无限的。我真的相信，动机非常重要，但是这并不意味着这样为别人做好事，不会给你本人带来好处，对不对？

像中大奖一样去实现自己的梦想，能帮助你创造更多的财富，因为创造财富的关键是释放你的创造力，找到一种方法能为别人做更多的事情。如果你找到一种方法能够为别人增加更多的价值，比现在其他人带来的增值都要多，那么你肯定也能找到办法，让自己的日子过得更好，让你自己个人的财富价值也积累得更多。这适用于你自己的生活，也适用于别人的生活。记住，我们在谈论如何加速你实现财务自由目标计划实施的时候，你如果想成为伟大的人，就要学会成为很多人的仆人。我们已经知道，生活偏爱那些帮助他人的人。通过帮助别人，你提升了自己，你对很多人慷慨，上天就会对你更慷慨，好运会主动来找你。

\begin{BookQuote}
给你自己内心的宁静。你值得快乐。你值得高兴。

\hfill ——汉娜·阿伦特
\end{BookQuote}

那么，你如何将钱放进你的梦想水桶？我给你讲三种方法。

第一种，当你一下子得到一大笔钱的时候，比如，一下子得到10000美元奖金，就像我们在上一章讲过的那个案例一样。

第二种，你的风险/成长水桶突然碰到利好消息，你的投资大涨，让你大赚。就像在拉斯维加斯赌城里赌博一样，现在也许是时候从桌子上拿走一些赢到手的钱，落袋为安了，这样你的风险就减小了。有个办法，我的很多学生都用过，就是把赚的钱拿走，分成几份，重新按照原来确定的资产配置来投资。比如，将1/3放进安全/安心水桶，1/3放进风险/成长水桶，1/3放进梦想水桶。比如，我前面说的那个获得10000美元奖金的案例，就是把3000美元拿出来放到你的梦想水桶里。

这样把属于你的风险/成长水桶里的钱拿出来1/3放进安全/安心水桶，就像从赌桌上把赢到手的钱拿出来一些来加快扩大你最安全的投资一样，安全/安心水桶里的钱越多，你的内心就会越安宁。继续留下1/3的钱放在风险/成长水桶里承担风险，去追求更大的潜在上行收益，但是你这样做用的是你赢的钱，你再亏也亏不了原来投资的本金。最后1/3的钱放到你的梦想水桶里，相当于你自己让自己中了一个大奖，你今天可以好好享受一下了。因为实现一些梦想会刺激你，激励你继续前进，继续用前面成功的方法行动，因为这让你想要赚更多的钱，存更多的钱，投资更加高效，因为你今天就享受到了回报，而不是等到将来的某天。

给你的梦想水桶注入更多钱的第三个方法是，从你的收入里拿出来固定比例的一小部分，找个地方攒起来。这样你越攒越多，慢慢地，这笔钱就够你去购买你梦想得到的东西了——它可能是你的第一套房子、一辆汽车、一次度假，也可能是那些很有趣的小玩意儿，让你今天内心感到快乐，外在变得美丽。但是，你一定要记住，这些钱都是你额外省出来的，你已经省出来的作为财务自由基金的钱，一分钱也不能动，财务自由是第一位的，先要自由，后要梦想！财务自由基金神圣不可侵犯。你总会找到方法，找到更多的钱，把更多的钱放进你的财务自由基金中，也把更多的钱放到你的梦想水桶里。这里我们快速回顾一下有关章节的内容：

\begin{itemize}
\item 存的更多。用多存的钱去投资更多
\item 赚的更多，用多赚的钱去投资更多
\item 减少费用和税收，用多省下来的钱去投资更多
\item 获得更高的投资收益率
\item 改变你的生活方式
\end{itemize}

这样做，你积蓄的钱就更多了，你就可以把这些储蓄拿出一部分投资，也拿出一部分让你的梦想成真——今天就实现梦想或者在不久后变为现实。

你用什么策略来装满这个梦想水桶？你会等着发奖金，或者股票大涨，或者你像我的朋友安吉拉那样从收入里拿出来一部分钱攒起来吗？一开始安吉拉觉得自己没有办法省钱，甚至为了达到财务自由也不行。

但是等到她看完本书，学会了7步简单投资法后，她发现，搬到佛罗里达州居住——因为那个州不征收个人收入所得税，这样她就能存起来个人收入的10\%作为她的财务自由基金，另外还能额外留出来8\%装进她的梦想水桶。原来要交的税，她现在不用交了，相当于税务局的人在帮她把钱装进梦想水桶里。这是不是很厉害？而且她还享受到了更好的天气！她仔细检查了自己的账户，找到了一个办法，使自己的钱变得更有税收效率，这样她又能从收入里拿出来2\%放进她的财务自由基金里。她总共能从收入里拿出来12\%存到她的财务自由基金里，除此之外，她还从收入里拿出8\%放进了梦想水桶。

要是你一开始告诉安吉拉她能够找到办法，把个人收入的20\%积攒起来，她肯定会说，你这家伙绝对是疯了。但是现在她不但能从个人收入里拿出12\%积攒起来，让自己未来的财务更加安全、更有保障，而且又额外从收入里拿出来8\%积攒起来，努力在短期内就实现自己的一些重要梦想，这些梦想让她非常兴奋——登喜马拉雅山，划着赛艇横渡大洋。安吉拉的大学专业是人类学，她一直梦想着能够追随著名的古生物学家路易斯·利基，到他在肯尼亚的研究所待上一段时间。她甚至已经拿到了邀请函，只不过她现在没有钱。但是只要她一直坚持按照她那个十分稳健的投资理财计划去做，持续不断地存钱，持续不断地投资，她不久就会有钱了。她既能够实现财产安全和财务独立，又能过上那种充满探索和冒险的生活，这多酷啊！还记得那个明天储蓄更多的策略吗？你可以决定，下一次你涨工资，也许可以拿出来3\%放到你的财务自由基金中，再拿出来1.5\%或者2\%放到你的梦想基金中——特别是如果有些梦想现在对于你来说非常重要的话，就像你需要储蓄一笔钱买第一套房子付首付一样，或者用于一个短期度假一样。竟然有这么多方式能够多存一些钱来实现你的梦想！

但是，我要告诉你一个秘密：最重要的事情是列出一个你的梦想清单。根据重要性、规模大小、时间长短，把你的梦想按照顺序排列起来。写下你必须要实现这些梦想的理由或者体验这些梦想的理由。我发现，如果你想找出来一个合适的收入储蓄百分比，却不知道你存这些钱的原因，那么你是存不下钱的。实现梦想的秘密在于，知道你真正想要的是什么，为什么你想要为这个梦想燃烧激情。突然间，你的创意就会释放出来，你会发现，一些新的办法让自己赚的更多，存的更多，给别人增加的价值更多，或者让税收效率变得更高，你成了更好的投资者，或者改变了你的生活方式——既改善了你的生活品质，又让你现在的一些梦想即刻实现，而不是等到未来。列出梦想清单，是实现所有这一切的关键所在。

但是，你今天就要做出决定！现在就花点儿时间，列出你的梦想清单吧。把你的梦想写下来，变成白纸黑字，好让这些梦想变得真实可见。你愿意为实现这些梦想储蓄多少钱？兴奋起来！行动起来！

\begin{BookQuote}
每一个伟大的梦想都开始于一个梦想者。

\hfill ——哈丽特·塔布曼
\end{BookQuote}

最终，你认为自己的总资产应该拿出来多大比例放到你的梦想水桶中？并不需要很多——也许是5\%或者10\%。但是请不要忘了好好犒劳一下自己。尽管努力让你的钱既安全又成长非常重要，但是绝对不要忘记寻找乐趣，去给予，活出完全的自我，就这样一边积累财富，一边享受生活，一路走向财务自由。因此，你既要有安全/安心水桶和风险/成长水桶，还要有梦想水桶。不要留着你的梦想水桶里的钱不花，直到将来为了应急才用。为什么不现在跑出去享受灿烂的阳光？享受生活，就在现在。

如果你做不到现在就享受生活，你可能最终会像那对夫妻一样可怜。这是一个朋友讲给我听的故事。这对老夫妻过日子精打细算，非常节省，一辈子都在省钱。后来，终于有一天，他们觉得钱攒得够多了，该存的钱都存够了，还有不少节余，足够他们进行一次加勒比海游轮之旅了。活了一辈子，怎么也得奢侈一把。他们的这次旅行，长达一周，他们坐在巨大的游轮上，绕着海岛旅行。你可以想象那个巨大游轮的样子：船上有游泳池，有攀岩设施，有几十个餐馆和几十个舞厅。这对夫妻玩得非常高兴，但是他们还是想省着点儿花，毕竟他们这钱攒得很不容易，辛辛苦苦地赚钱，精打细算地花钱，一分一厘地省钱，好不容易才攒够了退休养老的钱。他们可不想再额外花钱去吃大餐了。这趟旅行本身对于他们来说就已经够奢侈了。为了省钱。他们带了一个大箱子，里面装满了一盒又一盒的奶酪和饼干，用来在旅途上当饭吃，他们发誓坚决不吃那些价格昂贵的午餐和晚餐。

天气很好，这对夫妻在船上玩得也非常开心。但是每天一到午餐和晚餐时间，两人就不开心了。因为别的人都去吃自助餐了，巨大的餐桌，摆满了各种各样的食物，非常丰盛，龙虾、螃蟹、上等肋排、堆成小山的甜品，以及来自世界各地的上好葡萄酒。而他们俩只能回到房间里吃自己带的奶酪和饼干。他们不觉得难堪，也不觉得难过。他们很享受这趟一生也许只有一次的豪华旅行，他们很自豪自己过得很节俭，吃得很节省。但是，到了最后一天，他们实在受不了天天吃奶酪和饼干了，所以决定奢侈一把，到楼上吃一次丰盛的晚餐！他俩坐到其中一个巨大的自助餐桌旁边，大吃大喝，把他们的盘子装得满满的，里面都是他们这辈子吃过的最好的东西。

在吃了好多份甜点又喝了好多杯上等的葡萄酒后，他们实在吃不动了，于是叫来服务员买单。那个服务员一听“买单”很吃惊：“付什么钱？”这对夫妇说：“我们吃的这顿大餐啊。这些葡萄酒、甜点，我们吃的所有这些东西。”

服务员惊呆了，转过身来对这对夫妻说：“你不知道这趟旅程赠送一日三餐吗？”

旅程赠送一日三餐。你觉得这个比喻怎么样？所以，在你的这趟人生旅程里，不要只满足于奶酪和饼干。享受生命旅程一路赠送的所有东西吧。

另外再提醒一下：很多能让我们富有的东西都是免费的。还记得约翰·邓普顿爵士告诉过我们的吗？财富的秘密就是感恩。人生不只是我们取得的东西或者我们成就的东西，还有我们欣赏的东西。人生并不仅是一趟冒险之旅，人生也是我们应该花时间去享受的幸福之旅。你可以同时找到探险和快乐，在那些你爱的人身上，在你孩子那跳动的眼神里，在那些你爱的人快乐的脸上。你到处都可以中大奖，只要你能够醒悟，感受到你生活中的美好就在当下。所以不要发誓等到将来某天赚够了钱，现在就开始超越短缺。意识到你多么幸运，你拥有的财富竟然这么多：爱、欢乐、机会、健康、朋友、家人。不要只追求财富，现在就开始享受财富吧。

现在，我们已经学到了如何配置我们的投资，分散投资到不同类型、不同种类的资产上，把大类的资产分别装到两个不同的水桶——安全/安心水桶和风险/成长水桶。我们已经学到了，再拿出来另外一部分钱放进梦想水桶，让我们的人生过得更加有滋有味，也让我们更有动力为自己和别人做得更好。第四部分还有最后一章，我将教你三招，让你每年的投资收益率增加1\%\textasciitilde{}2\%，更重要的是确保你避免选时的错误。太多人都想预测市场行情的走势，以选择买卖时机，结果犯了大错。现在让我们从认知的力量中学习……
